\section{Introduction}
\label{sec:introduction}

Clusterized quantum magnets provide a direct route from the spectrum of a few-spin unit to the collective behavior of a quantum many-body system.
Their local multiplets can survive as low-energy degrees of freedom, become dispersive, or be removed from an experimental response by symmetry.
Clusterized chains therefore separate three questions that are often mixed together in a conventional phase diagram: which local state remains active, how neighboring units couple, and which resulting excitations are visible.

The recent spin-$1/2$ tetramer-chain study provides a useful reference because it treats magnetic order, entanglement structure, and excitation spectra within a common coupling diagram~\cite{li2026tetramer_verified}.
Replacing four spins by five is not, however, a minor geometric modification.
A unit cell with an odd number of half-integer spins retains a half-integer degree of freedom unless symmetry breaking or ground-state degeneracy intervenes.
Lieb--Schultz--Mattis-type results therefore constrain a unique symmetry-preserving gapped ground state under the relevant translation and spin-rotation assumptions~\cite{lieb1961,ogata2019}.
This constraint does not determine the phase diagram and does not by itself establish a globally gapless $c=1$ state.

Here we study a linear Heisenberg pentamer chain with two intra-cell exchanges, $J_1$ and $J_2$, and an inter-cell exchange $J_3$.
We set $J_1=1$ and define $\alpha=J_2/J_1$ and $\beta=J_3/J_1$.
The geometry differs from previously studied mixed-spin or differently connected pentamers, so their multiplet assignments cannot be transferred directly to the present model~\cite{haraldsen2011}.
Instead, we determine the complete isolated-cell spectrum, construct the low-energy projected chain, and test the resulting picture against static and dynamical many-body calculations.

Three results organize this work.
First, throughout the sampled positive-$\alpha$ region, the pentamer ground state is an even-parity spin-$1/2$ doublet and positive $\beta$ generates an antiferromagnetic effective chain.
Second, reflection parity makes the lowest excited doublet essentially dark at $q=\pi$ and concentrates the local inelastic weight in spin-$3/2$ multiplets.
Third, negative $\alpha$ produces a reproducible finite-size reconstruction of the dominant momentum.
The last result is not promoted to a thermodynamic phase boundary because two independently declared diagnostics do not localize the same adjacent interval.

This distinction between a resolved microscopic mechanism and a bounded phase-diagram statement is central to our analysis.
Production records are accepted only after numerical, transfer, checksum, and independent-validator gates.
We retain the raw spectral ringing, distinguish the Krylov dimension $k$ from the matrix-product-state bond dimension $\chi$, and report initial-state differences as systematic sensitivity rather than statistical uncertainty.

